\chapter{Bar--cobar and coefficient topology}
\label{v4-ch:bar-cobar-topology}

A nonempty ordered configuration of \(n\) positions has \(n-1\) gaps.
A word of positive geometric factors therefore has length at most
\(n\). This bound makes the cut bar and cobar finite at each
configuration. It does not bound the length of a resolution inside
one coefficient module. Such a resolution can be infinite even at
the empty configuration, where there are no geometric gaps.

We first prove the duality controlled by positive cuts. We then
construct the coefficient bar with its signs and augmentation,
and determine when it computes a derived tensor. The final sections
compare three ways of allowing infinite words. A series may
terminate at every finite precision while having unbounded length.
An unrestricted product permits still more series but can destroy
the continuity of an otherwise valid comparison.

\section{Positive words and finite duality}
\label{v4-sec:bar-positive}

Use the word category \(\mathcal W_X\), cut product
\(\star_{\mathrm{cut}}\), and empty unit \(K\) of
Chapter~\ref{v4-ch:words-collisions-units}. Write \(\star\) for this
product throughout the present section. For \(p\geq0\), put
\[
 \mathcal W_{\geq p}=\{F\in\mathcal W_X:F_n=0\text{ for }n<p\}.
\]
Thus \(\mathcal W_{\geq1}\) consists of the positive systems.

\begin{lemma}[Finite positive words]\label{v4-lem:bar-finite-words}
The cut product sends
\(\mathcal W_{\geq p}\times\mathcal W_{\geq q}\) to
\(\mathcal W_{\geq p+q}\). In particular,
\(F_1\star\cdots\star F_r\star M\) vanishes in arity \(n<r\)
when all \(F_i\) are positive, whether or not \(M\) is positive.
If \(V_r\in\mathcal W_{\geq r}\), then
\[
 \bigoplus_{r\geq0}V_r\ \xrightarrow{\sim}\ \prod_{r\geq0}V_r,
 \qquad
 \left(\bigoplus_{r\geq0}V_r\right)_n
       =\bigoplus_{r=0}^n(V_r)_n.
\]
\end{lemma}

\begin{proof}
A nonzero summand in a cut product assigns at least \(p\) positions
to its first factor and \(q\) to its second. The mixed assertion
uses one position for each positive factor and none necessarily
for \(M\).

The finite sums \(V_n=\bigoplus_{r\leq n}(V_r)_n\) form a word
system: for a monotone surjection \([n]\twoheadrightarrow[m]\),
exact extraordinary pullback takes this sum to
\(\bigoplus_{r\leq n}(V_r)_m\), whose terms with \(r>m\) vanish.
Since these pullbacks preserve colimits, \(V\) is the coproduct.
A family of maps \(T\to V_r\) gives, in each arity, a map to the
finite product defining \(V_n\). These maps commute with collisions
after every projection, and hence commute with them as maps to
\(V\). The same argument applies to coherent homotopies. It proves
the product universal property. No assertion about pullback of an
arbitrary infinite product is used.
\end{proof}

An augmented associative algebra in \(\mathcal W_X\) is an
associative algebra object \(A\) with an algebra map \(A\to K\).
It is \emph{connected} if this map is an equivalence at the empty
configuration. Its augmentation ideal
\(a=\operatorname{fib}(A\to K)\) is then positive, and the unit
splits \(A\simeq K\oplus a\). For a connected coaugmented coalgebra
\(C\), put \(c=\operatorname{cofib}(K\to C)\).
All algebra and coalgebra objects in this section retain their
higher associativity or coassociativity homotopies.

The \emph{cut bar} \(B_{\mathrm{cut}}A\) is the geometric
realization of the two-sided simplicial bar of \(K,A,K\).
Its faces multiply or augment, and its degeneracies insert the
unit. The \emph{cut cobar} \(\Omega_{\mathrm{cut}}C\) is the
totalization of the dual cosimplicial construction, whose cofaces
comultiply or coaugment and whose codegeneracies apply the counit.
For positive reduced parts their normalized terms are
\begin{equation}\label{v4-eq:bar-cut-normalized}
 \begin{split}
 B_{\mathrm{cut}}A&=K\oplus\bigoplus_{r\geq1}(a[1])^{\star r},\\
 \Omega_{\mathrm{cut}}C&=K\oplus\bigoplus_{r\geq1}(c[-1])^{\star r}.
 \end{split}
\end{equation}
These expressions include the normalized differentials and higher
attaching maps. They do not replace a homotopy-coherent algebra by
its graded object with zero differential.

\begin{theorem}[Finite cut duality]\label{v4-thm:bar-cut-duality}
Cut bar and cut cobar induce inverse equivalences between connected
augmented associative algebras and connected coaugmented
coassociative coalgebras in \((\mathcal W_X,\star,K)\).
Each normalized construction in arity \(n\) has at most \(n\)
positive factors. No global bound on word length or on coefficient
cohomological degrees is imposed.
\end{theorem}

\begin{proof}
The differential-module categories are presentable, and the
extraordinary pullbacks defining \(\mathcal W_X\) are accessible
and preserve colimits. Their limit is presentable. The associative
bar--cobar adjunction therefore supplies coherent maps
\[
 u_A:A\longrightarrow\Omega_{\mathrm{cut}}B_{\mathrm{cut}}A,
 \qquad
 v_C:B_{\mathrm{cut}}\Omega_{\mathrm{cut}}C\longrightarrow C.
\]
Here bar is left adjoint to totalized cobar. This adjunction uses
the full mapping spaces of algebra and coalgebra objects.\footnote{
J.~Lurie, \emph{Higher Algebra}, Theorem~5.2.2.17 and
Remark~5.2.2.19, pp.~827--828,
\href{https://people.math.harvard.edu/~lurie/papers/HA.pdf}{author's text}.
Warning~5.2.2.16 explains why no preservation of both arbitrary
realizations and totalizations by the tensor product is assumed.}

Splitting off the unit removes the degenerate simplicial terms;
the normalized term in degree \(r\) is \(a^{\star r}[r]\).
The normalized cosimplicial term is \(c^{\star r}[-r]\).
By Lemma~\ref{v4-lem:bar-finite-words}, the latter totalization
tower stabilizes at step \(n\) in arity \(n\). Its finite limits
commute with collision pullback. The resulting compatible family
is the totalization in \(\mathcal W_X\): maps into a word system
are the limit of its component mapping spaces, and that limit
commutes with the cosimplicial limit. This proves
\eqref{v4-eq:bar-cut-normalized} on the stated carrier.

Filter \(A\) increasingly by \(F_0A=K\), \(F_rA=A\) for
\(r\geq1\). In its associated graded, \(a\) has weight one and
zero multiplication. The induced filtration on
\(\Omega_{\mathrm{cut}}B_{\mathrm{cut}}A\) counts original
\(a\)-factors. At two factors, the single gap is either inside a
bar block or between two cobar factors. Splitting the block joins
these two copies of the coefficient by an identity map, with its
suspension sign. For \(r\) factors there are \(r-1\) such gaps.
The coefficient differential and block splitting preserve
\(r\); multiplication of original factors lowers it. The
associated graded is thus the cube that changes an inside gap to
an outside gap, with an identity map on the common coefficient at
each edge. The suspension signs give opposite signs around each
square. It is the tensor product of \(r-1\) two-term identity
complexes, tensored with the shifted coefficient
\(a^{\star r}\). For \(r>1\) it is contractible. The weights zero
and one are \(K\) and \(a\), and the associated graded map of
\(u_A\) is the identity on these terms.

For \(C\), use the decreasing filtration
\(F^0C=C\), \(F^1C=c\), \(F^2C=0\), where the splitting by the
counit regards \(c\) as a subobject. The associated graded reduced
coalgebra is primitive. In \(B_{\mathrm{cut}}\Omega_{\mathrm{cut}}C\),
the corresponding filtration counts original \(c\)-factors.
Cobar splitting increases this count, while bar multiplication
preserves it. On the associated graded the same gap cube runs from
outside bar cuts to inside product gaps. Its weights greater than
one are contractible; \(v_C\) is the identity on the surviving
weights zero and one.

After evaluation in arity \(n\), both filtrations are finite:
terms with more than \(n\) original positive factors vanish.
Their cofiber sequences prove that \(u_A\) and \(v_C\) are
equivalences there. Component evaluation is jointly conservative,
so they are equivalences of word systems. The construction in
filtered objects retains coherent maps. Higher products lower the
first filtration and higher coproducts raise the second; the same
associated graded calculation applies. The adjunction gives the
triangle identities and hence the asserted equivalence.
\end{proof}

\begin{corollary}\label{v4-cor:bar-vacuum-under}
For the vacuum coalgebra \(J\) of
Proposition~\ref{v4-prop:word-vacuum-coalgebra}, connected cut
coalgebras equipped with a coalgebra map \(J\to C\) correspond to
connected augmented cut algebras equipped with an algebra map
\(\Omega_{\mathrm{cut}}J\to A\).
\end{corollary}

\begin{proof}
An equivalence of categories induces an equivalence on the
categories of objects under any fixed object. Apply
Theorem~\ref{v4-thm:bar-cut-duality} to the object \(J\).
\end{proof}

Thus the local vacuum marking survives as an algebra map with
source \(\Omega_{\mathrm{cut}}J\). It is different from the
augmentation \(A\to K\). The theorem concerns the cut tensor on
words. A comparison with the full ordered atlas, or with prescribed
field operations, must preserve the additional geometry and maps
identified in Chapter~\ref{v4-ch:words-collisions-units}.

\section{The coefficient bar and its twisting maps}
\label{v4-sec:bar-coefficients}

Let \(R\) be a commutative ring and \(A\) a unital differential
graded \(R\)-algebra with a specified augmentation
\(\epsilon:A\to R\). Write \(\bar A=\ker\epsilon\), so
\(A=R\oplus\bar A\) as a complex. All tensor products in this
section are ordinary tensors over \(R\). The cohomological
differential has degree one. Put \(sV=V[1]\) and \(tV=V[-1]\);
the maps \(s\) and \(t\) have degrees minus one and one.

The \emph{coefficient bar} is
\begin{equation}\label{v4-eq:bar-coefficient}
 B_R A=T_R^c(s\bar A)
       =R\oplus\bigoplus_{r\geq1}(s\bar A)^{\otimes_R r}.
\end{equation}
The superscript \(c\) denotes the tensor coalgebra: its coproduct
deconcatenates each word at every position, including both ends.
The empty word gives the coaugmentation and counit. Its reduced
coproduct \(\bar\Delta\) omits the end cuts. A coaugmented
coalgebra is \emph{conilpotent} if every element of its reduced
part is killed by a sufficiently long iterated reduced coproduct.
This property is called \emph{conilpotence}.
Finite words make \(T_R^c(s\bar A)\) conilpotent.

Define its differential as the coderivation whose components with
one output factor are
\begin{equation}\label{v4-eq:bar-coefficient-signs}
 b_1(sa)=-s(d_Aa),\qquad
 b_2(sa,sb)=(-1)^{|a|}s(ab).
\end{equation}
Inserting either degree-one component after
\(w_1,\ldots,w_i\) has sign
\((-1)^{|w_1|+\cdots+|w_i|}\). On degree-zero algebra inputs,
the multiplication differential is therefore
\[
 b[a_1|\cdots|a_r]
    =\sum_{i=1}^{r-1}(-1)^{i-1}
                 [a_1|\cdots|a_ia_{i+1}|\cdots|a_r].
\]
The notation \([a_1|\cdots|a_r]\) abbreviates the tensor of the
suspended entries; its degree is \(\sum_i|a_i|-r\).

For a coaugmented counital differential graded coalgebra \(C\),
write \(\bar C=\ker(C\to R)\). Its direct-sum coefficient cobar is
the tensor algebra
\begin{equation}\label{v4-eq:bar-coefficient-cobar}
 \Omega_R C=T_R(t\bar C),\qquad
 d(tc)=-t(d_Cc)+
       \sum_{\bar\Delta c=\sum c'\otimes c''}
                    (-1)^{|c'|+1}(tc')(tc'').
\end{equation}
The displayed map extends by the graded Leibniz rule. The
coalgebra coproduct takes values in the ordinary algebraic tensor,
so it gives a finite sum on each element.

\begin{proposition}\label{v4-prop:bar-coefficient-identities}
The differentials in \eqref{v4-eq:bar-coefficient-signs} and
\eqref{v4-eq:bar-coefficient-cobar} square to zero. For conilpotent
\(C\), there are natural bijections
\[
 \operatorname{Hom}_{\mathrm{dgAlg}_R}(\Omega_R C,A)
 \cong\operatorname{Tw}_R(C,A)
 \cong\operatorname{Hom}_{\mathrm{dgCoalg}_R}(C,B_RA),
\]
where algebra maps preserve augmentations, coalgebra maps preserve
coaugmentations and counits, and \(\operatorname{Tw}_R(C,A)\)
consists of degree-one maps \(\tau:\bar C\to\bar A\) satisfying
\begin{equation}\label{v4-eq:bar-twisting}
 d_A\tau+\tau d_C+\mu_A(\tau\otimes\tau)\bar\Delta=0.
\end{equation}
The tensor of the two degree-one maps contributes
\((-1)^{|c'|}\) on \(c'\otimes c''\).
\end{proposition}

\begin{proof}
For the bar square, disjoint insertions cancel by their degree-one
signs. The overlapping multiplication terms on three inputs give
\((-1)^{|b|}s((ab)c-a(bc))\). Mixed terms vanish by the Leibniz
identity for \(d_A\), and the internal square is zero. These are
all possibilities. The corresponding cobar calculation uses that
\(\bar\Delta\) is a chain map and is coassociative: the two orders
of splitting one factor have opposite signs. Coderivation and
derivation compatibility hold by construction.

A map from the tensor algebra \(\Omega_RC\) is determined by
\(tc\mapsto\tau(c)\). Substitution in
\eqref{v4-eq:bar-coefficient-cobar} shows that its chain-map
condition is exactly \eqref{v4-eq:bar-twisting}. A coalgebra map
to the tensor coalgebra is determined by its length-one component
\(s\tau\); its length-\(r\) component must be
\[
 (s\tau)^{\otimes r}\bar\Delta^{(r-1)}.
\]
Conilpotence makes the sum of these components finite on each
input. Coassociativity proves the coalgebra identity. The
length-one part of its chain-map equation is
\eqref{v4-eq:bar-twisting}, with
\eqref{v4-eq:bar-coefficient-signs}; the coderivation identity
then proves the equation in every length. These constructions are
inverse and natural.
\end{proof}

This is a strict adjunction \(\Omega_R\dashv B_R\) for algebraic
tensor constructions and conilpotent coalgebras. It is distinct
from the general enhanced adjunction between geometric bar and
totalized cobar used in the finite cut theorem. On positive word
systems the finiteness theorem makes both comparisons equivalences.

\begin{proposition}[An algebra resolution]\label{v4-prop:bar-cobar-resolution}
For every augmented differential graded \(R\)-algebra \(A\),
the algebra map
\(\varepsilon_A:\Omega_RB_RA\to A\) that sends \(t(sa)\) to
\(a\), and a cobar generator represented by a longer bar word
to zero, is a quasi-isomorphism.
\end{proposition}

\begin{proof}
The desuspension \(B_RA\to A\) is a twisting map, so the preceding
proposition gives \(\varepsilon_A\). Filter by the number \(r\)
of original \(\bar A\)-entries. The target filtration is
\(F_0A=R\), \(F_rA=A\) for \(r\geq1\). The associated graded
source is the gap cube of the finite cut proof, now with
coefficient \(\bar A^{\otimes_R r}\). Each gap contributes a
two-term identity complex. Contracting the first gap contracts
their tensor product for \(r>1\), over any ring \(R\), without
a flatness hypothesis. In weights zero and one,
\(\varepsilon_A\) is the identity. Induction through the finite
filtration stages proves a quasi-isomorphism at each stage.
Every element belongs to some stage, and filtered unions of
\(R\)-modules are exact. Passing to that union proves the result.
No convergence of an infinite product is asserted.
\end{proof}

\section{Which derived tensor the bar computes}
\label{v4-sec:bar-derived-tensor}

A complex \(V\) of \(R\)-modules is \emph{\(K\)-flat} if
\(N\otimes_R V\) is acyclic for every acyclic complex \(N\).
This condition concerns the whole complex; for an unbounded
complex, degreewise flatness alone does not suffice.

\begin{proposition}\label{v4-prop:bar-derived-tensor}
If \(\bar A\) is \(K\)-flat over \(R\), then
\[
 B_RA\simeq R\otimes_A^{\mathbf L}R.
\]
The two modules \(R\) use the specified augmentation of \(A\).
\end{proposition}

\begin{proof}
Let \(\tau:B_RA\to A\) desuspend a word of length one and vanish
elsewhere. On \(L=A\otimes_R B_RA\), put
\[
 d_L(a\otimes c)=d_Aa\otimes c+(-1)^{|a|}a\otimes bc
      -(-1)^{|a|}\sum a\tau(c_{(1)})\otimes c_{(2)}.
\]
The twisting identity proves \(d_L^2=0\). Let
\(\epsilon_L:L\to R\) augment the length-zero term and vanish
on positive lengths. With \(\bar a=a-\epsilon(a)1\), define
\[
 h\bigl(a\otimes[ a_1|\cdots|a_r]\bigr)
       =-1\otimes[\bar a|a_1|\cdots|a_r].
\]
The internal terms cancel in \(d_Lh+hd_L\), as do contractions
away from the first slot. The first-slot term gives
\(\bar a\otimes c\); if \(r>0\), the terms that merge the
inserted slot with the next one leave \(\epsilon(a)\otimes c\).
Thus \(d_Lh+hd_L=1-\iota\epsilon_L\), where \(\iota\) includes
the empty word. This proves \(L\simeq R\).

The increasing length filtration has quotients
\(A\otimes_R(s\bar A)^{\otimes_R r}\). Tensoring one with an
acyclic right \(A\)-module \(N\) gives
\(N\otimes_R(s\bar A)^{\otimes_R r}\), which is acyclic by
\(K\)-flatness. The filtration is split on graded modules;
induction and exact filtered unions show that \(L\) is a
\(K\)-flat left \(A\)-resolution. Tensoring it with the right
augmentation module kills the first-slot term and leaves \(B_RA\).
\end{proof}

For a nonunital differential graded algebra \(P\), the tensor bar
on \(sP\) means the reduced bar of the unitization \(R\oplus P\),
augmented to \(R\). It need not be the reduced bar of a unital
algebra structure already present on \(P\). Indeed, if \(P\) has
a closed unit \(e\), inserting \(se\) at the beginning of a
nonempty word contracts this tensor bar onto \(R\):
\(b(se\otimes c)+se\otimes bc=c\). A reduced bar of \(P\)
itself instead requires a chosen augmentation and its kernel.
This distinction determines the meaning of the endpoint bars in
Chapter~\ref{v4-arith:gluing:chapter}.

\begin{example}[An infinite bar at the empty configuration]
\label{v4-ex:bar-empty-infinite}
Take \(A=k\oplus ka\), with \(|a|=1\), zero differential and
\(a^2=0\). Then \(v=sa\) has degree zero and
\[
 B_kA=\bigoplus_{r\geq0}kv^{\otimes r},\qquad b=0.
\]
Put this algebra at the empty component of \(\mathcal W_X\)
and zero at every positive arity. Its cut bar is exactly this
coefficient bar at the empty component. The augmentation is not
an equivalence there, so the algebra is outside the connected
class of Theorem~\ref{v4-thm:bar-cut-duality}. Arbitrarily long
coefficient words survive in the same cohomological degree.
\end{example}

\section{Length products and continuous Taylor maps}
\label{v4-sec:bar-length-completion}

For \(R\)-modules \(V_n\), the direct sum \(\bigoplus_nV_n\) contains
finite sums of words, whereas \(\prod_nV_n\) contains an arbitrary
word in every length. Give each \(V_n\) the discrete topology and
the product its product topology. A linear map of direct sums
has components \(F_{mn}:V_n\to W_m\).

\begin{proposition}[Continuity in length]\label{v4-prop:bar-row-finite}
A linear map \(F:\bigoplus_nV_n\to\bigoplus_mW_m\) extends to
a continuous map of the products of discrete modules if and only
if, for every \(m\), only finitely many \(F_{mn}\) are nonzero.
Such an extension is unique.

More generally, for a prime \(\ell\), let \(R=\mathbb Z_\ell\), give each \(V_n\) its
\(\ell\)-adic topology, and assume every \(W_m\) is separated
and complete for that topology. An \(R\)-linear \(F\) extends
continuously to the products if and only if, for each \(m,j\),
\[
 F_{mn}(V_n)\subseteq\ell^jW_m
       \quad\text{for all but finitely many }n.
\]
\end{proposition}

\begin{proof}
Continuity of the \(m\)-th coordinate at zero implies that its
kernel contains a basic product neighborhood. Outside finitely
many input coordinates that neighborhood contains each whole
module, so the corresponding \(F_{mn}\) vanish. Conversely,
finite row sums define a continuous map coordinate by coordinate.

For the adic statement, apply the same argument to the \(m\)-th
output modulo \(\ell^jW_m\). The asserted condition makes
\(\sum_nF_{mn}(x_n)\) converge in \(W_m\), uniformly modulo each
\(\ell^j\); the finite remaining terms are continuous because
they are \(R\)-linear. This constructs the extension. The direct
sum is dense in the product, and the target is separated, proving
uniqueness in both cases.
\end{proof}

A degree-zero coaugmentation-preserving map between tensor
coalgebras is specified by its Taylor components
\(f_r:V^{\otimes r}\to W\). Its value on a
word is the sum over consecutive subdivisions, applying one
\(f_r\) to each block. Every finite word has finitely many
subdivisions, but an output of length one can receive inputs of
every length. Proposition~\ref{v4-prop:bar-row-finite} detects
precisely this possible failure of continuity.

For example, let \(V=\mathbb Z v\) in degree zero and set
\(f_r(v^{\otimes r})=v\) for every \(r\geq1\). The resulting
coalgebra map is defined on finite words. The sums
\(z_N=\sum_{r=1}^Nv^{\otimes r}\) converge in the product of
discrete length spaces, but the length-one output is \(Nv\),
which does not converge. In the relative identity example of
Proposition~\ref{v4-arith:gluing:prop:noncontinuity}, these are
the Taylor components of the actual endpoint comparison.

A product also changes the recipient of deconcatenation. For
\(\widehat T(V)=\prod_{r\geq0}V^{\otimes r}\), the natural
recipient is
\[
 \prod_{p,q\geq0}(V^{\otimes p}\otimes V^{\otimes q}),
\]
with \((p,q)\)-component given by cutting length \(p+q\). It
need not be \(\widehat T(V)\otimes\widehat T(V)\).
Over a characteristic-zero field, take \(V=kv\) and the element
\(z=\sum_{r\geq0}r!v^{\otimes r}\). Its coproduct coefficient
matrix is \(((p+q)!)_{p,q\geq0}\). The leading \(d\)-square
minor has determinant
\[
 \det\bigl((p+q)!\bigr)_{0\leq p,q<d}
          =\prod_{i=0}^{d-1}(i!)^2\ne0.
\]
To see this, factor \(p!\) from row \(p\); the \(q\)-th column
is then a monic polynomial of degree \(q\) in \(p\), so its
determinant is the Vandermonde determinant. Every element of an
algebraic tensor is a finite sum of separated tensors and gives
a finite-rank coefficient matrix. Thus \(\Delta z\) cannot
belong to that algebraic tensor.

\section{Finite precision and derived inverse limits}
\label{v4-sec:bar-adic-completion}

Put \(R_j=\mathbb Z/\ell^j\), and let \(A_j\) be a compatible
tower of augmented differential graded \(R_j\)-algebras. Write
\(B_j=B_{R_j}A_j\). Define
\begin{equation}\label{v4-eq:bar-adic-completion}
 \widehat B=\varprojlim_jB_j,
 \qquad
 \widehat B\,\widehat\otimes\,\widehat B
       :=\varprojlim_j(B_j\otimes_{R_j}B_j).
\end{equation}
These are specified inverse limits, with discrete finite-precision
terms. No identification with another completed tensor is implicit.
The recipient of an iterated coproduct with \(r\) outputs is
\(\varprojlim_j B_j^{\otimes_{R_j}r}\).
An element has only finitely many nonzero lengths at each
precision, although that finite set can grow with \(j\).

\begin{proposition}\label{v4-prop:bar-adic-maps}
A compatible family of differential graded coalgebra maps
\(F_j:B_j\to B'_j\) induces a continuous map
\(\widehat F:\widehat B\to\widehat B'\) preserving the
coproduct with the recipients in
\eqref{v4-eq:bar-adic-completion}. If every \(F_j\) has a
compatible strict inverse, then \(\widehat F\) is a continuous
isomorphism.
\end{proposition}

\begin{proof}
Define the map by its projections. The chain and coproduct
identities hold after every projection and therefore in the
inverse limit. The output at precision \(j\) depends only on
the input at that precision, proving continuity. The compatible
inverses define its continuous inverse in the same way.
\end{proof}

For \(A_j=R_j\oplus R_ja\) with the square-zero structure of
Example~\ref{v4-ex:bar-empty-infinite}, the element
\[
 z=\sum_{r\geq1}\ell^rv^{\otimes r}
\]
belongs to \(\widehat B\). It is not elementwise conilpotent:
for any number of cuts, a sufficiently long word has a nonzero
coefficient in the corresponding tensor recipient. The Taylor
map with \(f_r(v^{\otimes r})=v\) is nevertheless continuous
on this restricted inverse limit. Its length-one value on \(z\)
is \(\ell(1-\ell)^{-1}v\). At precision \(j\) this is a finite
sum, so no analytic summation is part of its definition.

The preceding proposition is about strict maps of complexes. A
claim about their cohomology requires a different argument. For
a countable tower of complexes \(M_j\), with transition maps
\(q_j:M_{j+1}\to M_j\), define
\[
 R\!\varprojlim_jM_j
   =\operatorname{Cone}\left(
       \prod_jM_j\xrightarrow{1-T}\prod_jM_j\right)[-1],
 \qquad (Tx)_j=q_jx_{j+1}.
\]
This is the derived inverse limit of the tower of module complexes.

\begin{proposition}\label{v4-prop:bar-milnor}
If every \(q_j\) is surjective in each cohomological degree,
then the ordinary inverse limit represents this derived inverse
limit. For each \(q\) there is an exact sequence
\[
 0\longrightarrow\varprojlim{}^{1}_j H^{q-1}(M_j)
 \longrightarrow H^q(\varprojlim_jM_j)
 \longrightarrow\varprojlim_jH^q(M_j)\longrightarrow0,
\]
where \(\varprojlim{}^1\) is the cokernel of \(1-T\) on the
product of the indicated cohomology groups.
\end{proposition}

\begin{proof}
Surjectivity permits the recursive solution of
\(x_j-q_jx_{j+1}=y_j\), starting with any \(x_1\). Thus
\(1-T\) is degreewise surjective and its kernel is the ordinary
inverse limit. The resulting short exact sequence of complexes
identifies this kernel with the shifted cone. Products of modules
are exact, so their cohomology is the product of the cohomologies.
The long exact sequence of the cone gives the displayed sequence.
\end{proof}

In particular, commuting cohomology with the limit requires the
stated \(\varprojlim{}^1\) term to vanish. Scalar extension need
not commute with the limit either:
\[
 \mathbb Q_\ell\otimes_{\mathbb Z_\ell}
       \varprojlim_j\mathbb Z_\ell/\ell^j=\mathbb Q_\ell,
 \qquad
 \varprojlim_j\left(
       \mathbb Q_\ell\otimes_{\mathbb Z_\ell}
                    \mathbb Z_\ell/\ell^j\right)=0.
\]
The first identity uses \(\ell\)-adic completeness; in the
second, inverting \(\ell\) kills each finite quotient.

\section{Analytic bounds and the other bar constructions}
\label{v4-sec:bar-analytic-transverse}

There is a useful analytic counterpart to the continuity criterion.
For real or complex Banach spaces \(V,W\), use the completed
projective tensor powers \(V^{\widehat\otimes_\pi r}\). For
\(\rho>0\), let
\[
 T^1_\rho(V)=
 \left\{(x_r)_{r\geq0}:
          \sum_{r\geq0}\rho^r\|x_r\|_\pi<\infty\right\}.
\]
The displayed weighted sum is its norm, and the length-zero
summand is the scalar field. Suppose continuous
Taylor maps \(f_r:V^{\widehat\otimes_\pi r}\to W\) satisfy
\(\|f_r\|\leq ML^{r-1}\) for constants \(M,L>0\).

\begin{proposition}\label{v4-prop:bar-analytic-bound}
The subdivision map determined by these Taylor maps, acting as
the identity in length zero, extends to a bounded linear map
\(T^1_\rho(V)\to T^1_\sigma(W)\) whenever
\(L+\sigma M\leq\rho\). Its norm is at most one, and on the
positive-length part it is at most \(\sigma M/\rho\).
\end{proposition}

\begin{proof}
There are \(\binom{r-1}{m-1}\) consecutive subdivisions of a
length-\(r\) word into \(m\) nonempty blocks. Each corresponding
operator has norm at most \(M^mL^{r-m}\). Hence its entire
weighted output has norm at most
\[
 \sum_{m=1}^r\binom{r-1}{m-1}
       (\sigma M)^mL^{r-m}\|x_r\|_\pi
   =\sigma M(L+\sigma M)^{r-1}\|x_r\|_\pi
   \leq\frac{\sigma M}{\rho}\rho^r\|x_r\|_\pi.
\]
Sum over \(r\), retaining the scalar term separately. Finite
words are dense, so the bounded extension is unique.
\end{proof}

This estimate concerns the Taylor map. Continuity of the
differential and coproduct in the chosen recipients must also be
checked. For example, give the completed two-index tensor the norm
\(\sum_{p,q\geq0}\sigma^{p+q}\|z_{pq}\|_\pi\).
Deconcatenation from radius \(\rho\) to this recipient, for
\(\sigma<\rho\), is bounded by
\(\sup_{r\geq0}(r+1)(\sigma/\rho)^r\). At equal radii that
bound need not be finite. Algebraic identities on finite words
extend when all maps in those identities have continuous
extensions to a common domain and a separated target.

A \emph{transverse bar} uses another specified monoidal category
\((\mathcal E,\otimes_{\mathrm{tr}},\mathbf1_{\mathrm{tr}})\)
and an augmented associative algebra \(T\) for that tensor.
If geometric realizations exist and the tensor preserves them,
its two-sided bar is
\[
 B_{\mathrm{tr}}T
     =\mathbf1_{\mathrm{tr}}
          \otimes_T^{\mathbf L}\mathbf1_{\mathrm{tr}}.
\]
Its augmentation, coefficient category and convergence refer to
this tensor. The positive-cut bound applies to it only if an
additional compatible positive filtration is supplied.

Finally, for a covariant diagram \(F\) of complexes of \(k\)-vector
spaces on a small category \(\mathcal I\), its simplicial
replacement has terms
\[
 B_n(\mathcal I,F)
       =\bigoplus_{i_0\to\cdots\to i_n}F(i_0).
\]
Its direct-sum total complex has
\(\operatorname{Tot}^q=\bigoplus_{n\geq0}B_n^{q+n}\) and
differential \((-1)^n d_{\mathrm{int}}+\sum_i(-1)^id_i\), where
the \(d_i\) are the simplicial faces. It computes the derived colimit. In
Chapter~\ref{v4-arith:3d-acs-E1bar}, the ordered incidence construction
is of this kind. It needs an actual coefficient diagram on its
indexing category. Identifying that derived colimit with the
homology of a geometric or physical field theory is a further
comparison.
